\documentclass[11pt,a4paper]{article}
\usepackage[utf8]{inputenc}
\usepackage[T1]{fontenc}
\usepackage[ngerman]{babel}
\usepackage[margin=1in]{geometry}
\usepackage{xcolor}
\usepackage{hyperref}
\definecolor{primarycolor}{RGB}{33,37,41}
\definecolor{secondarycolor}{RGB}{0,102,204}
\hypersetup{colorlinks=true,linkcolor=secondarycolor,urlcolor=secondarycolor,citecolor=secondarycolor}
\setlength{\parindent}{0pt}
\setlength{\parskip}{0.8em}
\pagestyle{empty}
\begin{document}
{\large\bfseries\color{primarycolor} Katarzyna Wiśniewska}\\
Habsburgerstraße 41, 79104 Freiburg im Breisgau\\
k.wisniewska@example.com \textbar{} +49 761 555 0412

\vspace{0.6em}
Sozialstation Dreisamtal gGmbH

\vspace{0.6em}
\textbf{Bewerbung als Pflegehelferin für Ihren ambulanten Pflegedienst}

Sehr geehrte Damen und Herren,

seit April 2020 arbeite ich als Pflegehelferin in der stationären Altenpflege. Für die ausgeschriebene Stelle bringe ich langjährige Erfahrung in der Grundpflege, die Qualifikation Behandlungspflege LG1 und Deutschkenntnisse auf Niveau B2 mit. Ihr Leitsatz „Pflege mit Zeit und Respekt“ passt zu meinem ruhigen und geduldigen Umgang mit älteren Menschen und Menschen mit Demenz.

Im Seniorenzentrum Am Lindenhof übernehme ich die Grundpflege für 10 bis 12 Bewohnerinnen und Bewohner pro Schicht. Dazu gehören Körperpflege, Mobilisation, Lagerung und Hilfe beim Essen. In der Behandlungspflege nach LG1 führe ich Blutdruck- und Blutzuckermessungen durch, unterstütze bei Kompressionsstrümpfen und reiche Medikamente nach Plan. Früh-, Spät- und Wochenenddienste gehören zu meinem Arbeitsalltag. Die digitale Pflegedokumentation ist mir durch die Arbeit mit Vivendi PD vertraut; diese Erfahrung bietet eine Grundlage für die Einarbeitung in Ihr System MediFox.

Zuvor war ich als Krankenschwester auf einer internistischen Station in Polen tätig. Mein Anerkennungsverfahren als Pflegefachfrau beim Regierungspräsidium Stuttgart läuft. Für den Wechsel von der stationären in die ambulante Pflege bietet Ihre zweiwöchige begleitete Einarbeitung einen strukturierten Einstieg. Ihre festen Tourengruppen unterstützen zudem die Kontinuität im Umgang mit den betreuten Menschen, zu der ich meine zuverlässige und geduldige Arbeitsweise beitragen kann.

Gern erläutere ich Ihnen in einem persönlichen Gespräch, wie ich meine Pflegeerfahrung in Ihr Team einbringen kann. Ich freue mich über Ihre Einladung.

Mit freundlichen Grüßen

Katarzyna Wiśniewska
\end{document}